\section{Related Work}
\label{sec:related}

\subsection{Complexity of graph partitioning}

The minimum bisection problem --- partition a graph into two equal-weight
parts minimising the edge cut --- is NP-hard for general graphs
\citep{garey1976simplified}.
Garey, Johnson, and Stockmeyer~\citep{garey1976simplified} showed that
minimum bisection remains NP-hard even for degree-3 planar graphs via
reduction from 3-SAT.
Their reduction, however, does not impose connectivity on the two parts,
and hence does not directly cover the redistricting variant, which
requires each district to induce a connected subgraph.

Dyer and Frieze~\citep{dyer1985partitioning} proved that the problem of
partitioning a graph into $k$ \emph{connected} subgraphs of prescribed
sizes is NP-complete in general, and remains NP-complete for planar
graphs.
This is the result we use to anchor Theorem~\ref{thm:nphard}: the
redistricting problem requires connected parts by definition, so the
connected-partition variant is the correct source for the reduction.

The balanced $k$-partition problem (generalising bisection to $k$ parts,
without the connectivity requirement) is also
NP-hard~\citep{andreev2006balanced}, and the minimum edge-cost $k$-cut
problem is NP-hard for $k \geq 3$ even without balance
constraints~\citep{goldschmidt1994polynomial}.
The connectivity requirement of the redistricting problem is at least as
hard as the unconstrained case, since any algorithm for the unconstrained
problem that happened to produce connected parts would solve the
constrained problem.

\subsection{Approximation algorithms}

For general graphs, the best known approximation ratio for minimum bisection
is $O(\log n)$~\citep{andreev2006balanced}.
For planar graphs, the Arora-Rao-Vazirani (ARV)
framework~\citep{arora2009expander} achieves $O(\sqrt{\log n})$
approximation for Sparsest Cut via a semidefinite programming (SDP)
relaxation.
This bound applies to Sparsest Cut, not directly to Balanced $k$-Cut
with contiguity; extending it to balanced connected $k$-partition
with formal guarantees remains an open problem.

METIS~\citep{karypis1998} achieves near-optimal results empirically on
benchmark graphs but provides no formal approximation guarantee for
balanced connected partition beyond the planar separator
bound~\citep{lipton1979separator}.
The best known polynomial-time algorithm achieves no better than
$O(\log n)$ approximation for balanced $k$-partition in general; no
known algorithm achieves a tighter formal bound for the connected
variant.
KaHyPar~\citep{kahypar2017} provides stronger guarantees on hypergraph
instances but is not designed for planar adjacency graphs.

\subsection{Complexity of redistricting}

Discrete geometry approaches to redistricting~\citep{duchindist2020}
have characterised the solution space via Markov chains rather than
complexity theory.  Fifield et al.~\citep{fifield2020} proved mixing
bounds for redistricting Markov chains, showing that the solution space
is well-connected --- a result consistent with our seed sensitivity
analysis (B.7).

Computational redistricting has been characterised as an
optimisation problem by Mehrotra et al.~\citep{mehrotra1998optimal}.
Neither work provides complexity results for the minimum edge-cut variant.

\subsection{Runtime analysis of METIS}

Karypis and Kumar~\citep{karypis1998} proved $O(m \log n)$ runtime per
level for their multilevel scheme, where $m = |E|$ and $n = |V|$.
For planar graphs, $m = O(n)$ (by Euler's formula for planar graphs),
giving $O(n \log n)$ per bisection and $O(n \log k \log n)$ overall for
$\lceil \log_2 k \rceil$ levels of bisection.
Empirical scaling on non-planar graphs has been characterised for scientific
computing workloads~\citep{karypis1999multilevel}, but not for census
tract graphs with population-balance constraints.
